\documentclass[12pt, oneside, openany]{book}

\usepackage{mathptmx} % Contiene una fuente similar a Times New Roman

\usepackage[spanish, es-tabla]{babel} % Permite escritura en castellano
\usepackage[utf8]{inputenc} % Permite utilizar caracteres UTF8

\usepackage{graphicx} % Para la inclusión de gráficos e imágenes
\graphicspath{ {images/} } % Ruta para buscar las imágenes
\usepackage[a4paper,top=30mm,left=30mm,right=25mm,bottom=25mm,headheight=20mm]{geometry} % Configuración de los margenes de la página

% Paquetes para que funcione el formato.
\usepackage{titlesec}
\usepackage{setspace}
\usepackage{ragged2e}
\usepackage{fancyhdr}
\usepackage{lastpage}
\usepackage{stackengine}
\usepackage{array}
\usepackage[hidelinks]{hyperref}
\usepackage{enumitem}
\usepackage{float}
\usepackage{hypcap}
\usepackage{caption}
\usepackage{fancyvrb}

\usepackage{hyperref} % Paquete para que las referencias funcionen, y permite introducir links
\usepackage{xcolor} % Paquete para trabajar con colores (fondo de celdas, color del texto...)
% Se define un color gris desde su código RGB
\definecolor{gris}{RGB}{220,220,220}

\setcounter{secnumdepth}{3} % Para permitir numerar las sub-subsecciones

% Modifica el nombre de los índices al castellano
\addto\captionsspanish{
  \renewcommand{\contentsname}{Índice de contenido}
  \renewcommand{\listfigurename}{Índice de figuras}
  \renewcommand{\listtablename}{Índice de tablas}
}

% Formateo de los nombres de los apartados:
\titleformat{\chapter}[block]
  {\normalfont\Huge\bfseries\singlespacing}{\thechapter.}{1em}{\Huge}
\titlespacing*{\chapter}{0pt}{-62pt}{0pt}

\titleformat{\section}[block]
  {\normalfont\Large\bfseries}{\thesection.}{4pt}{\Large}
\titlespacing*{\section}{0pt}{\baselineskip}{0pt}

\titleformat{\subsection}[block]
  {\normalfont\large\bfseries}{\thesubsection.}{4pt}{\normalsize\large}
\titlespacing*{\subsection}{0pt}{0pt}{0pt}

\titleformat{\subsubsection}[block]
  {\normalfont\normalsize\bfseries}{\thesubsubsection.}{4pt}{\normalsize}
\titlespacing*{\subsubsection}{0pt}{0pt}{0pt}

\def\tablename{Tabla}

%% Variables para portada y cabeceras
%% Cambiar los valores para cada documento!!!
\def\title{Segunda entrega de prácticas}
\def\subject{Pruebas y Despliegue de Software}
\def\author{Juan Francisco Mier Montoto}
\def\target{www.mier.info}
\def\authorid{UO283319 -~PL1}
\def\date{febrero 2024}
\def\org{Escuela Politécnica de Ingeniería de Gijón}
\def\area{Grado en Ingeniería Informática en Tecnologías de la Información}

\def\ORG{\expandafter\MakeUppercase\expandafter{\org}}
\def\AREA{\expandafter\MakeUppercase\expandafter{\area}}
\def\SUBJECT{\expandafter\MakeUppercase\expandafter{\subject}}

\def\OK{OK\cellcolor{green!25}}
\newcommand{\MAL}[1]{MAL ({#1})\cellcolor{red!25}}
\captionsetup{justification=centering}
\setlength{\headheight}{65pt}

\fancyhf{}
\fancyhead[L]{\includegraphics[height=16mm]{style/square.png}
  \hspace{1em} \Longstack[l] {
    \textbf{\SUBJECT} \newline
    \textbf{\title}}
  \newline \leftmark{}
}
\fancyhead[R]{\bfseries{Hoja \hyperlink{toc}{\thepage}~de~\pageref{LastPage}}}
\fancyfoot[C]{\href{\target}{\author}}
\renewcommand{\headrulewidth}{0pt} % default is 0pt
\renewcommand{\footrulewidth}{0.4pt} % default is 0

\fancypagestyle{plain}{%
  \fancyhf{}
  \fancyhead[L]{\includegraphics[height=16mm]{style/square.png}
    \hspace{1em} \Longstack[l]{
      \textbf{\SUBJECT} \newline
      \textbf{\title}}}
  \fancyhead[R]{\bfseries{Hoja \hyperlink{toc}{\thepage}~de~\pageref{LastPage}}}
  \fancyfoot[C]{\href{\target}{\author}}
  \renewcommand{\headrulewidth}{0pt} % default is 0pt
  \renewcommand{\footrulewidth}{0.4pt} % default is 0pt
}

\pagestyle{fancy}

\restylefloat{table}



\begin{document}

\rmfamily % Fuente tipo Romana

% Portada de la memoria
\begin{titlepage}
    \centering
    \bfseries {
        \null{}
        \vspace{0cm}
        \begin{table}[h]
            \centering
            \begin{tabular}{m{10cm} m{1cm} m{3cm}}
                \vspace{0.2cm}
                \includegraphics[width=86mm]{style/full.png} &  & \vspace{1.52mm} \includegraphics[width=23mm]{style/square.png} \\
            \end{tabular}
        \end{table}

        \vspace{3\baselineskip}

        \Large{\ORG{} \\ \vspace{3\baselineskip}}
        \large {
            \AREA{} \\ \vspace{3\baselineskip}
            \subject{} \\ \vspace{2\baselineskip}

            % TRABAJO FIN DE GRADO/MÁSTER Nº XXXXXXXXX \vspace{\baselineskip} \\
            \title{} \\ \vspace{3\baselineskip}

            \author{} \\
            \authorid{} \\
            % TUTOR/ES: \\
            % D. APELLIDO1 APELLIDO2, Nombre \\
            % D. APELLIDO1 APELLIDO2, Nombre \\  \vspace{\baselineskip}

            \vspace{2\baselineskip}
            FECHA:\@ \date{}
        }
    }
\end{titlepage}


% Índice de contenido
\addcontentsline{toc}{chapter}{Índice de contenido} % Añade la referencia al índice de contenido
\hypertarget{toc}{}
\tableofcontents
\newpage

\justify{} % Texto justificado
\setlength{\parskip}{\baselineskip} % Separación entre párrafos de 1 linea
\onehalfspacing{}

%% El contenido de la memoria, dividido en capítulos:
\chapter{Introducción}


\chapter{Especificación de las pruebas}

\section{Clases de equivalencia}
A partir de la especificación de requisitos, se han identificado las siguientes clases de equivalencia:

\subsection{Entrada}
\begin{itemize}
	\item \textbf{Fecha de inscripción}
		\begin{enumerate}
			\item Fecha anterior a Inicio 1~(\textcolor{red}{I})
			\item Fecha entre Inicio 1 e Inicio 2
			\item Fecha entre Inicio 2 y cierre
			\item Fecha entre cierre y celebración de la actividad
			\item Fecha posterior a la celebración de la actividad~(\textcolor{red}{I})
		\end{enumerate}
	\item \textbf{Género}
		\begin{enumerate}
			\item Hombre
			\item Mujer
			\item Otro~(\textcolor{red}{I})
		\end{enumerate}
	\item \textbf{Edad}
		\begin{itemize}
			\item Menor que 18~(\textcolor{red}{I})
			\item Menor que 35
			\item Menor que 40
			\item Menor que 45
			\item Menor que 50
			\item Menor que 55
			\item Menor que 60
			\item Menor que 65
			\item Menor que 70
			\item Mayor que 70
			% \item Número muy grande (ej. 1000)~(\textcolor{red}{I})
		\end{itemize}
	\item \textbf{¿Inscripción ya existente?}
			\begin{itemize}
				\item Sí~(\textcolor{red}{I})
				\item No
			\end{itemize}
	\item \textbf{¿Quedan plazas vacías?}
			\begin{itemize}
				\item Sí
				\item No~(\textcolor{red}{I})
			\end{itemize}
\end{itemize}

\subsection{Salida}
\begin{itemize}
	\item \textbf{Categoría}
		\begin{itemize}
			\item Masculina: SM, M-35, M-40, M-45, M-50, M-55, M-60, M-65, M-70, M-70
			\item Femenina: SF, F-40, F-45, F-50, F-55, F-60, F-65, F-70, F-70
		\end{itemize}
	\item \textbf{Cuota}
		\begin{itemize}
			\item Cuota reducida 1
			\item Cuota reducida 2 (opcional)
			\item Cuota de inscripción
		\end{itemize}
	\item \textbf{Resultado}
		\begin{itemize}
			\item Inscripción aceptada
			\item Inscripción rechazada~(\textcolor{red}{I})
		\end{itemize}
\end{itemize}

\section{Combinatoria}
A la hora de determinar las situaciones a probar, se han generado \textit{combinaciones} de
distintos casos de equivalencia anteirormente descritos para reducir el número de pruebas a
realizar teniendo en cuenta la funcionalidad y los requisitos descritos.

Puesto que la combinación de las entradas de \textit{edad} y \textit{sexo} resultan en la
asignación a diferentes grupos dentro de una carrera, resulta lógico combinar ambas clases
para reducir el número de pruebas a realizar. La técnica combinatoria elegida en este caso
es \textit{pair-wise}, que consiste en combinar cada clase de equivalencia con todas las demás
clases de equivalencia al menos una vez. Se escoge esta técnica porque se considera importante
probar todas las combinaciones posibles de las clases de equivalencia, pero no es necesario
probar todas las combinaciones de todas las clases de equivalencia.

\begin{table}[ht]
	\centering
	\rowcolors{2}{white}{gray!15}
	\begin{tabular}{|c|c|c|}
		\hline
		\textbf{Edad} & \textbf{Masculino} & \textbf{Femenino} \\
		\hline
		\hline
		$<18$ & Inválido \cellcolor{red!25} & Inválido \cellcolor{red!25} \\
		$<35$ & SM & SF \\
		$<40$ & M-35 & SF \\
		$<45$ & M-40 & F-40 \\
		$<50$ & M-45 & F-45 \\
		$<55$ & M-50 & F-50 \\
		$<60$ & M-55 & F-55 \\
		$<65$ & M-60 & F-60 \\
		$<70$ & M-65 & F-70 \\
		$\geq 70$ & M-70 & F-70 \\
		\hline
	\end{tabular}
\end{table}

\section{AVL~(valores límite)}
Para obtener una mayor granularidad en la especificación de las pruebas, se han identificado
los valores límite de las clases de equivalencia. Estos valores límite se utilizan para
determinar los casos de prueba que se deben realizar para garantizar que el sistema se
comporta correctamente en los límites de las clases de equivalencia.

Obviamente, no se aplica AVL a TODAS las clases de equivalencia, sino a aquellas que se
consideran más relevantes.

\begin{itemize}
	\item \textbf{Fecha de inscripción}: se aplica AVL para comprobar que en los límites de las
		fechas de inscripción, el sistema se comporta correctamente. |TODO| completar
	\item \textbf{Edad}: se aplica AVL en el límite entre mayor y menor de edad, y en algunos
		límites de las categorías de edad para comprobar que todo funciona correctamente, es
		decir: \begin{itemize}
			\item Atleta menor de edad por un día en el día de la carrera
			\item Atleta con mayoría de edad cumplida el mismo día de la carrera
			\item Atleta con 35 años cumplidos el día de la carrera
		\end{itemize}
\end{itemize}

\input{sections/3_diseño.tex}
\chapter{Reporte de fallos}
Después de probar el sistema con el diseño y los datos anteriormente mencionados,
se detectan los siguientes fallos (en orden de relevancia):

\begin{itemize}
	\item No se contempla la categoría M-35, los corredores que deberían pertenecer a
		esta categoría se inscriben en la categoría SM. (SPI1.8)
	\item No se calcula bien el límite de mayoría de edad en los casos probados, ya que no se
		permite inscribir a atletas de 18 años cumplidos el día de la carrera. (SPI1.3)
	\item Los organizadores pueden crear competiciones con fechas de inicio, cierre y día de carrera
		en cualquier orden, siendo posible que la fecha de cierre sea posterior al día del evento, o
		que no exista plazo lógico de inscripción (entre otras combinaciones). (SC3, SC4, SC5, SC6)
	\item Se admiten inscripciones cuando el día de inscripción está entre el día de la carrera y
		tres días antes, cuando no debería ser posible. (SPI2.6)
	\item Se admiten atletas sin país. (SVI3)
\end{itemize}

\chapter{Resumen ejecutivo}
Pese a que se han detectado los fallos anteriormente mencionados, se ha de destacar la complejidad
de la lógica del negocio. La cantidad de situaciones que se han de probar es muy elevada, y la
combinatoria de las mismas es muy compleja.

Gracias al diseño de las pruebas, se han podido detectar algunos fallos importantes, sobre todo aquellos
relacionados con la lógica de negocio (dejando de lado la validación de datos, que se considera más trivial
tanto de probar como de solventar).


\end{document}
